% ===================== 5.1.2 建模与求解 =====================
% 模型建立（三段式，占最大篇幅）：
%   第1段（引文）：针对问题X，本节将从XX角度建立XX模型...具体过程如下。
%   第2段（展开）：按 \textbf{步骤N：标题} 编号展开，每步 6-7 个公式，
%     公式间用文字串联（"将式(X)代入式(Y)"），每式后紧跟变量说明段
%     （"其中，$X$表示...，$Y$表示..."，自然段不分点）。
%   第3段（桥接）：综上所述，本节完成了XX模型的建立，接下来进行求解。
% 四类问题的编号步骤骨架（详见 references/paper-writing.md）：
%   优化类：  1.目标函数的建立 -> 2.约束条件的建立 -> 3.核心机理建模 -> 4.求解算法建模
%   评价类：  1.评价指标体系构建 -> 2.指标权重确定（熵权/AHP 完整推导） -> 3.综合评价方法
%   预测/分类：1.特征工程 -> 2.模型选定与构建 -> 3.模型训练与评估
%   机理类：  1.机理分析框架 -> 2.关键变量关系建模 -> 3.模型整合与参数确定
% 模型求解：引文 -> 表/图 -> 分析。结果数值用 longtable，趋势对比用图（图宽 0.8\textwidth）。

\subsubsection{模型的建立}

\textbf{步骤1：}...描述。
\begin{equation}
  f(x) = ...
\end{equation}
其中，$x$表示...，$f(x)$表示...。

\textbf{步骤2：}...描述。
\begin{equation}
  g(x) = ...
\end{equation}
其中，...。

\subsubsection{模型的求解}

...求解过程与结果分析，关键数值以 longtable 呈现，并引用每张图表。
